\section{Automated Attack Prompt Construction} \label{sec:automated_prompts}
The manual construction of attack prompts for prompt injection attacks, as shown in~\cref{sec:manual_prompts}, is able to achieve high attack success rates, mainly by applying general prompting practices to adversarial settings.
While this approach is well-suited for early experimentation and simple red-teaming tasks, it is also time-consuming and difficult to apply at scale.
Automating parts of manual prompt construction workflows like the one proposed by \perez alleviates these problems to a certain extent.
In the following, we are presenting and discussing the partial automation of prompt injection attacks using the example of a semi-automated framework proposed by \deng~\cite{}.

\subsection{Semi-Automated Attack Framework}
In their attack framework, \deng automate two steps that were carried out manually in the approach of Perez and Ribeiro.
The first automation is the rephrasing of attack prompts.
As pointed out in~\cref{sec:basic_prompting_techniques}, \perez adapt the attack prompts by manually applying generally applicable prompting techniques.
In contrast to that, \deng leverage the text generation capabilities of LLMs to produce new attack prompts based on an existing set of older attack prompts~\cite{}.
To do so, they use the existing attack prompts as in-context learning examples and task the so-called attack LLM with generating new attack prompt proposals.\\

\begin{figure}[htb]
    \centering
    \includegraphics[width=\linewidth]{add_later.png}
    \caption{}
    \Description[<>]{<>}
    \label{fig:semi_automated_framework}
\end{figure}

The second step that \deng automate is the evaluation of proposed attack prompts.
While \perez use comparisons of actual and expected model responses to measure the effectiveness of injected prompts, \deng leverage the text understanding capabilities of LLMs~\cite{}.
They use an LLM-as-a-judge approach by tasking a so-called evaluation LLM with estimating the harmfulness of the attacked target LLM's response.
The evaluation LLM is tasked with assigning a harmfulness score on a scale of 1 to 10 to each of the target LLM's responses~\cite{}.
The most effective new attack prompts are then added to the attack prompt dataset, and the process can be repeated.\\

All steps of the attack framework proposed by \deng are depicted in~\cref{fig:semi_automated_framework}.
By iteratively repeating the automated prompt generation and evaluation steps, \deng can create a large set of successful attack prompts based on a small initial attack prompt set without any time-intensive human labor~\cite{}.
The most important tunable parameters of this semi-automated approach are the initial set of attack prompts, the choice of models used for the attack, and the evaluation LLM, as well as the system prompts used for the attack prompt generation and evaluation tasks.

\subsection{Automation Impact}
The most important advantage of the semi-automated attack framework proposed by \deng is that it eliminates most of the manual steps required by previous approaches and is therefore easily scalable.
This becomes evident by the fact that \deng are able to produce novel semi-automatic attack prompt (SAP) datasets with hundreds of unique prompts~\cite{}.
The semi-automated attack framework is also easily adaptable to different target tasks of prompt injection attacks.
While \deng use the target task of outputting harmful content in various topics to create the SAP datasets, the target task can be changed by adapting the system prompts of the attack and evaluation LLM accordingly~\cite{}.
\\

Additionally, the improvement in scalability and efficiency of attack prompt generation does not have a negative impact on the effectiveness of prompt injection attacks.
\deng prove this by comparing their SAP datasets with the manually created Dual-Use and the automatically generated BAD+ datasets.
They find that the SAP30 dataset outperforms both the Dual-Use and the BAD+ datasets on a total of four different variants of the GPT-3.5 and Alpaca-LoRA model families~\cite{}.
A selection of the evaluation results measured by \deng can be seen in~\cref{tab:deng_model_comparison}.
It is, however, important to notice that \deng apply the same harmfulness score used to generate the SAP datasets as a measure of comparison to evaluate SAP against the other datasets.
Since the compared datasets were created without taking this specific harmfulness score into consideration, an external metric could provide a more neutral evaluation.\\

\begin{table}[h]
    \centering
    \caption{Harmfulness scores assigned to outputs from prompt injection attacks on different models and datasets as measured by \deng The exact model identifiers of the GPT-3.5 models are \texttt{gpt3.5-turbo-0301} and \texttt{text-davinci-003}~\cite{}.}
    \label{tab:deng_model_comparison}
    \begin{tabular}{@{}lcccccc@{}}
        \toprule
         & \multicolumn{4}{c}{\textbf{Model}} \\
        \cmidrule{2-5}
        \textbf{Dataset} & \multicolumn{2}{c}{GPT-3.5} & \multicolumn{2}{c}{Alpaca-LoRA} \\
        \cmidrule(r){2-3} \cmidrule(l){4-5}
        & turbo & davinci & 7B & 13B \\
        \midrule
        \textbf{Dual-Use} & 5.41 & 6.35 & 6.66 & 6.33 \\
        \textbf{BAD+} & 0.63 & 1.87 & 4.12 & 3.44 \\
        \textbf{SAP30} & \textbf{8.70} & \textbf{7.18} & \textbf{8.80} & \textbf{8.72} \\ 
        \bottomrule
    \end{tabular}
\end{table}

An important limitation of the automation framework proposed by \deng is that it requires an initial attack prompt set as a baseline to start the LLM-based rephrasing process.
If no suitable baseline dataset is available for a specific attacker goal task, a manual dataset has to be constructed~\cite{}.
This adds time-intensive and expensive human labor to the framework, preventing the semi-automated approach from being fully automated.
In addition to the increasing effort, the baseline attack prompt dataset might also have an impact on the framework's performance.
The effectiveness of the new attack prompts, which are created by the framework through LLM-based rephrasing, is determined by the type of model, the model's system prompt, and the provided in-context learning examples.
While \deng do not investigate the impact of using different baseline attack prompt datasets for the same attack task, the choice of baseline prompts likely impacts the framework's performance.
For example, an LLM is unlikely to come up with completely new attack prompt strategies rather than rephrasing attack approaches that are included in the baseline attack prompts.

\subsection{Supervised Fine-Tuning and RLHF}
Combining the attack prompt effectiveness measurements from \perez and \deng allows us to draw some qualitative conclusions about the robustness of different fine-tuning approaches.
As mentioned in~\cref{sec:RLHF_instr_following}, \perez find that the \texttt{text-davinci-002} model from the GPT-3.5 series is significantly more susceptible to prompt injection attacks than the \texttt{text-davinci-001} model from the InstructGPT series~\cite{}.
The main difference between the two models is that \texttt{text-davinci-002} is fine-tuned using a supervised fine-tuning approach, while \texttt{text-davinci-001} is fine-tuned with RLHF~\cite{}.\\

\deng make a similar observation in their comparison of \texttt{text-davinci-003} from the GPT-3.5 series with models from the Alpaca series.
Based on the harmfulness scores from~\cref{tab:deng_model_comparison}, they conclude that the GPT-3.5 based models are more robust against prompt injection attacks than the Alpaca models, which are fine-tuned using a self-instruct framework~\cite{}.
This conclusion aligns with the findings by \perez and suggests that fine-tuning with RLHF leads to increased model robustness compared to supervised fine-tuning.\\

It is important to mention several aspects about this qualitative conclusion.
First of all, the observation by \deng that GPT-3.5 series models are more robust against prompt injection attacks than Alpaca models is based on the comparison of models that likely differ in size.
While the parameter numbers of the Alpaca models used by \deng are known, the parameter numbers of the GPT-3.5 series remain undisclosed~\cite{}.
At least the \texttt{text-davinci-003} model is likely to have a much higher parameter number compared to the Alpaca models, which could impact robustness.
Second, there is no data available to directly compare the findings by \perez and \deng
The comparison is only based on the findings that in both works, RLHF fine-tuned models are less susceptible to prompt injection attacks than comparable models using supervised fine-tuning.
Lastly, the comparison of RLHF and supervised fine-tuning relies on the empirical comparison of existing models.
The fact that models fine-tuned with RLHF are more robust does not mean that RLHF is in principle less susceptible to prompt injection attacks or that supervised fine-tuning is unable to produce sufficiently robust LLMs.